\chapter{Tres en raya paso a paso}

\section{Estado inicial del ejemplo}

Usaremos una posición intermedia para que el árbol sea suficientemente pequeño como para dibujarlo, pero suficientemente rico como para mostrar decisiones reales.

\[
S_0=\board{X}{O}{X}{\blank}{O}{\blank}{\blank}{\blank}{\blank}
\]

Es el turno de MAX, que juega con $X$. Las casillas vacías son
\[
\{4,6,7,8,9\}.
\]
Por tanto, el factor de ramificación local es
\[
 b(S_0)=5.
\]

\begin{stepbox}{Paso 1: identificar movimientos legales}
Desde $S_0$, MAX puede jugar en cualquiera de las cinco casillas vacías. Nombramos los movimientos por la casilla elegida:
\[
 a_4,a_6,a_7,a_8,a_9.
\]
Cada movimiento genera un hijo distinto.
\end{stepbox}

\section{Primer nivel del árbol}

\[
\begin{forest}
game tree
[$S_0$\\$X$ mueve
  [$S_4$\\$X$ en 4]
  [$S_6$\\$X$ en 6]
  [$S_7$\\$X$ en 7]
  [$S_8$\\$X$ en 8]
  [$S_9$\\$X$ en 9]
]
\end{forest}
\]

\section{Tableros de los hijos}

\begin{center}
\begin{tabular}{cc}
$S_4=\board{X}{O}{X}{X}{O}{\blank}{\blank}{\blank}{\blank}$ &
$S_6=\board{X}{O}{X}{\blank}{O}{X}{\blank}{\blank}{\blank}$ \\\[5mm]
$S_7=\board{X}{O}{X}{\blank}{O}{\blank}{X}{\blank}{\blank}$ &
$S_8=\board{X}{O}{X}{\blank}{O}{\blank}{\blank}{X}{\blank}$ \\\[5mm]
\multicolumn{2}{c}{$S_9=\board{X}{O}{X}{\blank}{O}{\blank}{\blank}{\blank}{X}$}
\end{tabular}
\end{center}

\begin{stepbox}{Paso 2: revisar si algún hijo es terminal}
Un estado es terminal si alguien ganó o si no quedan movimientos. Ninguno de los cinco hijos anteriores contiene tres $X$ alineadas ni tres $O$ alineadas, y todos conservan casillas vacías. Por tanto, ninguno es terminal.
\end{stepbox}

\section{Segundo nivel: responde MIN}

Ahora MIN juega con $O$. Analicemos primero el hijo $S_4$.

\[
S_4=\board{X}{O}{X}{X}{O}{\blank}{\blank}{\blank}{\blank}
\]

Las casillas vacías son $\{6,7,8,9\}$. Por tanto, MIN tiene cuatro respuestas.

\[
\begin{forest}
game tree
[$S_4$\\$O$ mueve
  [$S_{4,6}$\\$O$ en 6]
  [$S_{4,7}$\\$O$ en 7]
  [$S_{4,8}$\\$O$ en 8]
  [$S_{4,9}$\\$O$ en 9]
]
\end{forest}
\]

Los tableros son:

\begin{center}
\begin{tabular}{cc}
$S_{4,6}=\board{X}{O}{X}{X}{O}{O}{\blank}{\blank}{\blank}$ &
$S_{4,7}=\board{X}{O}{X}{X}{O}{\blank}{O}{\blank}{\blank}$ \\\[5mm]
$S_{4,8}=\board{X}{O}{X}{X}{O}{\blank}{\blank}{O}{\blank}$ &
$S_{4,9}=\board{X}{O}{X}{X}{O}{\blank}{\blank}{\blank}{O}$
\end{tabular}
\end{center}

\begin{stepbox}{Paso 3: detectar amenazas inmediatas}
En $S_{4,8}$, MIN coloca $O$ en la casilla 8. Entonces la columna central queda
\[
2,5,8 = O,O,O.
\]
Por tanto, $S_{4,8}$ es terminal y MAX pierde:
\[
\Utility(S_{4,8})=-1.
\]
\end{stepbox}

\section{Retropropagación Minimax local}

Como $S_4$ es un nodo de MIN, su valor es el mínimo de los valores de sus hijos:
\[
V(S_4)=\min\{V(S_{4,6}),V(S_{4,7}),V(S_{4,8}),V(S_{4,9})\}.
\]
Ya sabemos que
\[
V(S_{4,8})=-1.
\]
Como en tres en raya los valores posibles son $-1,0,1$, ningún valor puede ser menor que $-1$. Por tanto,
\[
V(S_4)=-1.
\]

\begin{intuitionbox}
MAX no debe jugar en la casilla 4, porque MIN responde en la casilla 8 y gana inmediatamente formando la columna central.
\end{intuitionbox}

\section{Lectura Alpha--Beta del mismo fragmento}

Al comenzar en la raíz:
\[
\alpha=-\infty,\qquad \beta=+\infty.
\]
Supongamos que el primer hijo examinado es $S_4$. Entramos al nodo $S_4$, donde mueve MIN. MIN busca un hijo con valor lo más bajo posible.

Cuando MIN encuentra $S_{4,8}$ con valor $-1$, sabe que puede forzar como máximo $-1$ para MAX desde $S_4$. Como $-1$ es el peor valor posible para MAX, no necesita seguir buscando respuestas mejores para MIN. Localmente, el subárbol queda resuelto.

\begin{stepbox}{Paso 4: qué se aprendió}
El análisis de $S_4$ no requiere expandir todos sus descendientes. Basta encontrar una refutación inmediata. Esta es la intuición operativa de la poda: una rama se descarta cuando ya no puede mejorar la decisión del jugador que la consideraba.
\end{stepbox}

\section{Árbol anotado parcial}

\[
\begin{forest}
game tree
[$S_0$\\MAX\\$\alpha=-\infty,\beta=+\infty$
  [$S_4$\\MIN
    [$S_{4,6}$]
    [$S_{4,7}$]
    [$S_{4,8}$\\terminal\\$-1$, fill=softRed]
    [$S_{4,9}$]
  ]
  [$S_6$]
  [$S_7$]
  [$S_8$]
  [$S_9$]
]
\end{forest}
\]

\section{Explicación del nodo marcado}

El nodo $S_{4,8}$ representa la secuencia:
\[
X\text{ juega }4,\qquad O\text{ responde }8.
\]
El tablero resultante es
\[
\board{X}{O}{X}{X}{O}{\blank}{\blank}{O}{\blank}.
\]
La columna central contiene $O$ en las posiciones $2,5,8$. Por tanto, MIN gana. Como la utilidad se mide desde MAX,
\[
\Utility(S_{4,8})=-1.
\]

\section{Continuación esperada del capítulo}

Las siguientes entregas desarrollarán los otros hijos $S_6,S_7,S_8,S_9$ con el mismo nivel de detalle:

\begin{enumerate}
\item expansión completa del subárbol relevante;
\item detección de amenazas;
\item evaluación de hojas;
\item retropropagación Minimax;
\item comparación con Alpha--Beta;
\item árbol coloreado con ramas podadas;
\item tabla de evolución de $\alpha$ y $\beta$.
\end{enumerate}

\begin{summarybox}
Este capítulo no pretende resolver todo tres en raya en una página. El objetivo es mostrar cómo se construye el árbol correctamente: estado por estado, movimiento por movimiento, valor por valor.
\end{summarybox}
